\documentclass[../../main.tex]{subfiles}
\begin{document}
\section{流出}\label{section:ch09-流出}

流为流形上的许多几何构造提供了技术工具。这些构造大多基于如下一般定理，它描述流在某些子流形附近的行为。

\begin{theorem}\label{theorem:ch09-20}[流出定理]
假设 $M$ 是光滑流形，$S\subseteq M$ 是嵌入的 $k$ 维子流形，$V\in\mathfrak{X}(M)$ 是处处不与 $S$ 相切的光滑向量场。设 $\theta\colon\mathscr{D}\to M$ 是 $V$ 的流，令 $\mathcal{O}=(\mathbb{R}\times S)\cap\mathscr{D}$，$\Phi=\theta|_{\mathcal{O}}$。
\begin{enumerate}[(a)]
\item $\Phi\colon\mathcal{O}\to M$ 是浸入。
\item $\partial/\partial t\in\mathfrak{X}(\mathcal{O})$ 与 $V$ 是 $\Phi$-相关的。
\item 存在光滑正函数 $\delta\colon S\to\mathbb{R}$ 使得 $\Phi$ 在 $\mathcal{O}_\delta$ 上的限制是单射，其中 $\mathcal{O}_\delta\subseteq\mathcal{O}$ 是流域
\begin{align}
\mathcal{O}_\delta=\{(t;p)\in\mathcal{O}:|t|<\delta(p)\}. \label{align:ch09-09}
\end{align}
因此，$\Phi(\mathcal{O}_\delta)$ 是 $M$ 的包含 $S$ 的浸入子流形，且 $V$ 与该子流形相切。
\item 若 $S$ 的余维为 $1$，则 $\Phi|_{\mathcal{O}_\delta}$ 是到 $M$ 的开子流形的微分同胚。
\end{enumerate}
\end{theorem}

\begin{remark}
子流形 $\Phi(\mathcal{O}_\delta)\subseteq M$ 称为从 $S$ 沿 $V$ 的\textbf{流出}。
\end{remark}

\begin{figure}[ht]
\centering
\includegraphics[width=0.8\textwidth]{fig-09-7.png}
\caption{流形流出}\label{figure:ch09-flowout}
\end{figure}

\begin{proof}
首先证明 (b)。固定 $p\in S$，令 $\sigma\colon\mathscr{D}(p)\to\mathbb{R}\times S$ 为曲线 $\sigma(t)=(t;p)$。则 $\Phi\circ\sigma(t)=\theta(t;p)$ 是 $V$ 的积分曲线，故对任意 $t_0\in\mathscr{D}(p)$ 有
\begin{align}
\mathrm{d}\Phi_{(t_0;p)}\left(\frac{\partial}{\partial t}\bigg|_{(t_0;p)}\right)=(\Phi\circ\sigma)'(t_0)=V_{\Phi(t_0;p)}. \label{align:ch09-10}
\end{align}

接下来证明 (a)。$\Phi$ 在 $\{0\}\times S$ 上的限制是微分同胚 $\{0\}\times S\leftrightarrow S$ 与嵌入 $S\hookrightarrow M$ 的复合，故是嵌入。因此 $\mathrm{d}\Phi_{(0;p)}$ 在 $T_pS$（视为 $T_{(0;p)}\mathcal{O}\cong T_0\mathbb{R}\oplus T_pS$ 的子空间）上的限制是包含 $T_pS\hookrightarrow T_pM$。若 $(E_1,\dots,E_k)$ 是 $T_pS$ 的任意基，则 $\mathrm{d}\Phi_{(0;p)}$ 将 $T_{(0;p)}\mathcal{O}$ 的基
\begin{align*}
\left(\partial/\partial t\big|_{(0;p)},E_1,\dots,E_k\right)
\end{align*}
映为 $(V_p,E_1,\dots,E_k)$。由于 $V_p$ 不与 $S$ 相切，这个 $(k+1)$ 元组线性无关，故 $\mathrm{d}\Phi_{(0;p)}$ 单射。

为证明 $\mathrm{d}\Phi$ 在其他点单射，我们如等变秩定理的证明那样论证。给定 $(t_0;p_0)\in\mathcal{O}$，令 $\tau_{t_0}\colon\mathcal{O}\to\mathbb{R}\times S$ 为平移 $\tau_{t_0}(t;p)=(t+t_0;p)$。由 $\theta$ 的群律，下图交换（其中水平映射可能仅在包含 $(0;p_0)$ 和 $p_0$ 的开子集上有定义）：
\begin{center}
\begin{tikzcd}[ampersand replacement=\&]
\mathcal{O} \ar[r,"\tau_{t_0}"] \ar[d,"\Phi"'] \& \mathcal{O} \ar[d,"\Phi"] \\
M \ar[r,"\theta_{t_0}"'] \& M.
\end{tikzcd}
\end{center}
上图中两个水平映射都是局部微分同胚。取微分，得
\begin{center}
\begin{tikzcd}[ampersand replacement=\&]
T_{(0;p_0)}\mathcal{O} \ar[r,"\mathrm{d}(\tau_{t_0})_{(0;p_0)}"] \ar[d,"\mathrm{d}\Phi_{(0;p_0)}"'] \& T_{(t_0;p_0)}\mathcal{O} \ar[d,"\mathrm{d}\Phi_{(t_0;p_0)}"] \\
T_{p_0}M \ar[r,"\mathrm{d}(\theta_{t_0})_{p_0}"'] \& T_{\Phi(t_0;p_0)}M.
\end{tikzcd}
\end{center}
由于水平映射是同构，两个竖直映射有相同的秩。由于已证 $\mathrm{d}\Phi_{(0;p_0)}$ 满秩，$\mathrm{d}\Phi_{(t_0;p_0)}$ 亦然。这完成了 $\Phi$ 是浸入的证明。

接下来证明 (c)。给定点 $p_0\in S$，选取 $M$ 中以 $p_0$ 为中心的关于 $S$ 的切片坐标卡 $(U,(x^i))$，使得 $U\cap S$ 是 $x^{k+1}=\cdots=x^n=0$ 的集合（其中 $n=\dim M$）。由于 $V$ 不与 $S$ 相切，$V_{p_0}$ 的最后 $n-k$ 个分量之一，比如 $V^j(p_0)$，必非零。必要时缩小 $U$，可假设存在常数 $c>0$ 使得
\begin{align}
|V^j(p)|\geqslant c\quad\text{对所有 }p\in U. \label{align:ch09-11}
\end{align}
由于 $\Phi^{-1}(U)$ 在 $\mathbb{R}\times S$ 中开，可取 $\varepsilon_{p_0}>0$ 和 $p_0$ 在 $S$ 中的邻域 $W_{p_0}$ 使得 $[-\varepsilon_{p_0},\varepsilon_{p_0}]\times W_{p_0}\subseteq\mathcal{O}$ 且 $\Phi([-\varepsilon_{p_0},\varepsilon_{p_0}]\times W_{p_0})\subseteq U$。在这些局部坐标中将 $\Phi$ 的分量函数写为 $\Phi(t;p)=(\Phi^1(t;p),\dots,\Phi^n(t;p))$。由于 $\Phi$ 是流的限制，分量函数 $\Phi^j$ 满足
\begin{align*}
\frac{\partial\Phi^j}{\partial t}(t;p)=V^j\bigl(\Phi(t;p)\bigr),\qquad V^j(0;p)=0.
\end{align*}
由\eqref{align:ch09-11}和微积分基本定理，$|\Phi^j(t;p)|\geqslant c|t|$，因而对 $(t;p)\in(-\varepsilon_{p_0},\varepsilon_{p_0})\times W_{p_0}$，推出 $\Phi(t;p)\in S$ 当且仅当 $t=0$。

选取从属于 $S$ 的开覆盖 $\{W_p:p\in S\}$ 的光滑单位分解 $\{\varphi_p:p\in S\}$，定义 $f\colon S\to\mathbb{R}$ 为
\begin{align}
f(q)=\sum_{p\in S}\varepsilon_p\,\varphi_p(q). \label{align:ch09-12}
\end{align}
则 $f$ 光滑且为正。对每个 $q\in S$，只有有限个 $p\in S$ 使得 $\varphi_p(q)>0$；若 $p_0$ 是这些点中使 $\varepsilon_{p_0}$ 最大者，则
\begin{align*}
f(q)\geqslant\varepsilon_{p_0}\sum_{p\in S}\varphi_p(q)=\varepsilon_{p_0}.
\end{align*}
由此，若 $(t;q)\in\mathcal{O}$ 使得 $|t|<f(q)$，则 $(t;q)\in(-\varepsilon_{p_0},\varepsilon_{p_0})\times W_{p_0}$，故 $\Phi(t;q)\in S$ 当且仅当 $t=0$。

令 $\delta=\frac12 f$。我们将证明 $\Phi|_{\mathcal{O}_\delta}$ 是单射，其中 $\mathcal{O}_\delta$ 由\eqref{align:ch09-09}定义。假设对某 $(t;q),(t_0;q_0)\in\mathcal{O}_\delta$ 有 $\Phi(t;q)=\Phi(t_0;q_0)$。必要时重新命名点，可设 $f(q_0)\leqslant f(q)$。我们的假设意味着 $\theta_t(q)=\theta_{t_0}(q_0)$，$\theta$ 的群律继而推出 $\theta_{t-t_0}(q)=q_0\in S$。$(t;q)$ 和 $(t_0;q_0)$ 属于 $\mathcal{O}_\delta$ 这一事实推出
\begin{align*}
|t-t_0|\leqslant|t|+|t_0|<\tfrac12 f(q)+\tfrac12 f(q_0)\leqslant f(q),
\end{align*}
这迫使 $t=t_0$，从而 $q=q_0$。

只剩 (d)。若 $S$ 的余维为 $1$，则 $\Phi|_{\mathcal{O}_\delta}$ 是同维流形之间的单射光滑浸入，故是嵌入（\cref{proposition:ch04-22}(d)）且是到开子流形的微分同胚（\cref{proposition:ch05-01}）。
\end{proof}

\subsection*{正则点与奇点}

若 $V$ 是 $M$ 上的向量场，点 $p\in M$ 称为 $V$ 的\textbf{奇点}，若 $V_p=0$；否则称为\textbf{正则点}。下一个命题表明从正则点和奇点出发的积分曲线行为非常不同。

\begin{proposition}\label{proposition:ch09-21}
设 $V$ 是光滑流形 $M$ 上的光滑向量场，$\theta\colon\mathscr{D}\to M$ 是 $V$ 生成的流。若 $p\in M$ 是 $V$ 的奇点，则 $\mathscr{D}(p)=\mathbb{R}$ 且 $\theta^{(p)}$ 是常值曲线 $\theta^{(p)}(t)\equiv p$。若 $p$ 是正则点，则 $\theta^{(p)}\colon\mathscr{D}(p)\to M$ 是光滑浸入。
\end{proposition}

\begin{proof}
若 $V_p=0$，则由 $\gamma(t)\equiv p$ 给出的常值曲线 $\gamma\colon\mathbb{R}\to M$ 显然是 $V$ 的积分曲线，故由唯一性和极大性它必等于 $\theta^{(p)}$。

为验证第二个陈述，我们证明其逆否命题：若 $\theta^{(p)}$ 不是浸入，则 $p$ 是奇点。$\theta^{(p)}$ 不是浸入的假设意味着存在某 $s\in\mathscr{D}(p)$ 使得 $\theta^{(p)\prime}(s)=0$。记 $q=\theta^{(p)}(s)$。则上一段的论证推出 $\mathscr{D}(q)=\mathbb{R}$ 且对所有 $t\in\mathbb{R}$，$\theta^{(q)}(t)=q$。由\cref{theorem:ch09-12}(b) 推出 $\mathscr{D}(p)=\mathbb{R}$ 也成立，且对所有 $t\in\mathbb{R}$，群律给出
\begin{align*}
\theta^{(p)}(t)=\theta_t(p)=\theta_{t-s}\bigl(\theta_s(p)\bigr)=\theta_{t-s}(q)=q.
\end{align*}
令 $t=0$ 得 $p=q$，故 $\theta^{(p)}(t)\equiv p$ 且 $V_p=\theta^{(p)\prime}(0)=0$。
\end{proof}

若 $\theta\colon\mathscr{D}\to M$ 是流，点 $p\in M$ 称为 $\theta$ 的\textbf{平衡点}，若对所有 $t\in\mathscr{D}(p)$，$\theta(t;p)=p$。\cref{proposition:ch09-21}表明光滑流的平衡点恰为其无穷小生成元的奇点。

下一个定理完全描述了在微分同胚意义下，向量场在正则点附近的样子。

\begin{theorem}\label{theorem:ch09-22}[正则点附近的标准型]
设 $V$ 是光滑流形 $M$ 上的光滑向量场，$p\in M$ 是 $V$ 的正则点。存在 $p$ 某邻域上的光滑坐标 $(s^i)$，使得 $V$ 在其中有坐标表示 $\partial/\partial s^1$。若 $S\subseteq M$ 是任意嵌入超曲面，$p\in S$ 且 $V_p\notin T_pS$，则坐标还可取得使 $s^1$ 是 $S$ 的局部定义函数。
\end{theorem}

\begin{proof}
若未给出超曲面 $S$，选取以 $p$ 为中心的任意光滑坐标 $(U,(x^i))$，令 $S\subseteq U$ 为由 $x^j=0$ 定义的超曲面，其中 $j$ 选得使 $V^j(p)\neq 0$。（回忆 $p$ 是 $V$ 的正则点。）

无论 $S$ 是给定的还是如上构造的，由于 $V_p\notin T_pS$，必要时缩小 $S$ 可使 $V$ 处处不与 $S$ 相切。流出定理继而说存在流域 $\mathcal{O}_\delta\subseteq\mathbb{R}\times S$，使得 $V$ 的流限制为从 $\mathcal{O}_\delta$ 到包含 $S$ 的开子集 $W\subseteq M$ 的微分同胚 $\Phi$。$\mathcal{O}_\delta$ 中有 $(0;p)$ 的乘积邻域 $(-\varepsilon,\varepsilon)\times W_0$。选取光滑局部参数化 $X\colon\Omega\to S$，其像包含于 $W_0$，其中 $\Omega$ 是 $\mathbb{R}^{n-1}$ 的开子集，坐标记为 $(s^2,\dots,s^n)$。于是由
\begin{align*}
\Psi(t;s^2,\dots,s^n)=\Phi\bigl(t,X(s^2,\dots,s^n)\bigr)
\end{align*}
给出的映射 $\Psi\colon(-\varepsilon,\varepsilon)\times\Omega\to M$ 是到 $M$ 中 $p$ 的邻域的微分同胚。由于微分同胚 $(t;s^1,\dots,s^n)\mapsto(t,X(s^2,\dots,s^n))$ 将 $\partial/\partial t$ 前推到自身，而 $\Phi_*(\partial/\partial t)=V$，故 $\Psi_*(\partial/\partial t)=V$。因此 $\Psi^{-1}$ 是光滑坐标卡，$V$ 在其中有坐标表示 $\partial/\partial t$。将 $t$ 重命名为 $s^1$ 即完成证明。
\end{proof}

标准型定理的证明实际上提供了一种寻找使给定向量场 $V$ 化为标准型的坐标的技巧，至少在相应 ODE 组可显式求解时如此：从 $V$ 不与之相切的超曲面 $S$ 和局部参数化 $X\colon\Omega\to S$ 开始，构造复合映射 $\Psi(t;s)=\theta_t(X(s))$，其中 $\theta$ 是 $V$ 的流。所求坐标映射即为 $\Psi$ 的逆。该过程最好用例子说明。

\begin{example}\label{example:ch09-23}
设 $W=x\partial/\partial y-y\partial/\partial x$ 是 $\mathbb{R}^2$ 上的向量场。我们在\cref{example:ch09-08}(b) 中计算了 $W$ 的流。点 $(1;0)\in\mathbb{R}^2$ 是 $W$ 的正则点，因为 $W_{(1;0)}=\partial/\partial y|_{(1;0)}\neq 0$。由于 $W$ 在该处有非零的 $y$ 坐标，可取 $S$ 为 $x$ 轴，由 $X(s)=(s;0)$ 参数化。定义 $\Psi\colon\mathbb{R}^2\to\mathbb{R}^2$ 为
\begin{align*}
\Psi(t;s)=\theta_t(s;0)=(s\cos t,s\sin t),
\end{align*}
然后在 $(1;0)$ 邻域内局部用 $(x;y)$ 解出 $(t;s)$，得到 $(1;0)$ 邻域内的坐标映射
\begin{align}
(t;s)=\Psi^{-1}(x;y)=\bigl(\tan^{-1}(y/x),\sqrt{x^2+y^2}\bigr). \label{align:ch09-13}
\end{align}
容易验证在这些坐标下 $W=\partial/\partial t$。（正如你可能注意到的，它们只是换了名字的极坐标。）
\end{example}

标准型定理表明，正则点附近的流在微分同胚意义下恰如 $\mathbb{R}^n$ 中沿平行坐标直线的平移。因此，流的所有有趣局部行为都集中在其平衡点附近。平衡点周围的流可表现出多种多样的行为，例如包围平衡点的闭轨道、当 $t\to+\infty$ 或 $-\infty$ 时收敛到平衡点的轨道，以及更复杂的现象。一些典型的二维例子见图 9.8。平面中平衡点附近流的局部行为的系统研究可在许多 ODE 教科书中找到，如 [BD09]。流形上流的全局和长时间行为的研究，称为\textbf{光滑动力系统理论}，是一个深刻的学科，在数学内外都有许多应用。

\begin{figure}[ht]
\centering
\includegraphics[width=0.8\textwidth]{fig-09-8.png}
\caption{平衡点附近的流的例子}\label{figure:ch09-flows-equilibrium}
\end{figure}

\end{document}